% Generated draft: requires mathematical and source review.
\documentclass[11pt]{article}
\usepackage[margin=25mm]{geometry}
\usepackage{fontspec}
\setmainfont{Latin Modern Roman}
\usepackage{amsmath,amssymb,mathtools,unicode-math}
\setmathfont{Latin Modern Math}
\usepackage{enumitem,needspace,xurl,hyperref}
\setlength{\parindent}{0pt}
\setlength{\parskip}{5pt}
\setlength{\emergencystretch}{3em}
\newcommand{\N}{\mathbb{N}}
\newcommand{\Z}{\mathbb{Z}}
\newcommand{\Q}{\mathbb{Q}}
\newcommand{\R}{\mathbb{R}}
\newcommand{\C}{\mathbb{C}}
\newcommand{\F}{\mathbb{F}}
\newcommand{\A}{\mathbb{A}}
\newcommand{\PP}{\mathbb{P}}
\newcommand{\E}{\mathbb{E}}
\newcommand{\eps}{\varepsilon}
\newcommand{\dd}{\,\mathrm{d}}
\begin{document}
\section*{UnsolvedMath ID: 30005795}
\textbf{Source problem code:} catalogue entry
\subsection{Definitions and mathematical statement}
Let \(X\) be a smooth Fano threefold and let \(P\in S\subset X\). Here \(P\) denotes a point, whereas \(P(u)\) below denotes a divisor. The supplied record does not specify the ground field, the precise hypotheses on \(S\), or the definition of \(\tau\). For each real parameter \(u\in[0,\tau]\), the record writes a Zariski decomposition
\[
-K_X-uS=P(u)+N(u),
\]
with positive part \(P(u)\) and effective negative part \(N(u)\). The intended notion of Zariski decomposition on this threefold, and the hypotheses ensuring its existence throughout the interval, require clarification. Write \(\operatorname{Supp}N(u)\) for the union of the prime divisors occurring with nonzero coefficient in \(N(u)\), and write \((-K_X)^3\) for the anticanonical triple intersection number.

\textbf{Reconstructed target in the supplied registry.} Assuming
\[
\exists u\in[0,\tau]\quad P\in\operatorname{Supp}N(u),
\]
the registry asks for a constant \(K>0\), uniform in \(F\), such that
\[
\frac{3}{(-K_X)^3}
\int_0^\tau
\bigl(P(u)^2\cdot S\bigr)
\operatorname{ord}_F\!\bigl(N(u)|_S\bigr)\,\dd u
\leq K A_S(F)
\qquad
\text{for every prime divisor }F\text{ over }X.
\]
This is a transcription of the reconstructed target, with unresolved qualifications, rather than an unambiguous mathematical conjecture. The permitted dependence of \(K\) on \(X,S,P\) and the decomposition data is not specified.

\textbf{Unresolved valuation domain.} A prime divisor over a variety is a prime divisor on a normal variety equipped with a proper birational morphism to that variety. The reconstructed statement quantifies over prime divisors over \(X\), but the supplied original excerpt instead says ``for every prime divisor \(F\) over \(S\).'' The terms \(A_S(F)\) and \(\operatorname{ord}_F(N(u)|_S)\) naturally require valuations over \(S\); they are not intrinsically defined for an arbitrary prime divisor over \(X\).

For clarity of notation only, if \(S\) is normal and \(\mathbb Q\)-Gorenstein and \(F\) is a prime divisor on a normal birational model \(f:T\to S\), the usual log discrepancy is
\[
A_S(F)=1+\operatorname{coeff}_F(K_T-f^*K_S).
\]
If \(N(u)|_S\) is a defined \(\mathbb R\)-Cartier divisor, its order along such an \(F\) is
\[
\operatorname{ord}_F(N(u)|_S)
=\operatorname{coeff}_F\!\bigl(f^*(N(u)|_S)\bigr).
\]
These conditional definitions do not supply missing hypotheses or resolve the conflict between the two quantifiers. In particular, the record must clarify the nature and singularities of \(S\), whether the restriction \(N(u)|_S\) exists, the meaning and integrability of the integrand, and whether the allowed divisors must satisfy a center condition involving \(P\). No corrected choice of valuation domain or additional hypotheses is asserted here.
\subsection{Short English statement}
For permanent UnsolvedMath ID \(30005795\), the supplied question seeks a bound, uniform in the allowed divisorial valuations, for an integral measuring the negative part of \(-K_X-uS\) along \(S\), in the case where the point \(P\) lies in that negative part for some \(u\). The reconstructed statement uses divisors over \(X\), while the original excerpt uses divisors over \(S\). This conflict and missing geometric hypotheses prevent a definitive formulation. The supplied open status is dated curation, not an independently verified conclusion.
\subsection{Sources}
\begin{itemize}
\item Supplied local UnsolvedMath registry record, permanent ID \(30005795\), title: \emph{Valuation Bounds When the Center Meets the Negative Part}; exact source problem code: \texttt{catalogue entry}. Collection: Oberwolfach Reports Open Problems; proposed year: 2024. No individual proposer is named.
\item \emph{K-Stability, Birational Geometry and Mirror Symmetry} (2024), Oberwolfach Reports, pp.~827--870, DOI: \texttt{10.4171/owr/2024/14}. Supplied source URL and literature-source link: \url{https://doi.org/10.4171/owr/2024/14}. The supplied background describes these reports as records of mathematical workshops and their open-problem discussions.
\item The supplied original excerpt asks what happens when \(P\in\operatorname{Supp}N(u)\) for some \(u\in[0,\tau]\), includes the phrase ``for every prime divisor \(F\) over \(S\),'' and contains a corrupted surrounding formula concerning the complementary case. The reconstructed registry statement instead quantifies over divisors over \(X\). The registry labels its statement \texttt{corrected\_verified} and says it was reconstructed and checked against the official report; this is supplied verification metadata, not an independent verification in this draft.
\item Literature assessment checked on \textbf{2026-08-22}: the supplied curation labels the problem open, reports that the 2024 report poses the negative-part case as an unresolved technical question seeking a uniform bound, and reports that no later proof resolving that exact case was located. This is supplied dated curation, not independently verified. The supplied background expressly requests independent verification of time-sensitive status and cited claims.
\end{itemize}
\end{document}
